\documentclass[11pt,a4paper]{article}

\usepackage[utf8]{inputenc}
\usepackage[T1]{fontenc}
\usepackage{lmodern}
\usepackage{amsmath,amssymb}
\usepackage{booktabs}
\usepackage{tabularx}
\usepackage{longtable}
\usepackage[margin=2.5cm]{geometry}
\usepackage{enumitem}
\usepackage{caption}
\usepackage{float}          % for [H] placement
\usepackage{parskip}
\usepackage{titlesec}
\usepackage{tcolorbox}
\usepackage[table,dvipsnames]{xcolor}

\usepackage{hyperref}
\hypersetup{
    colorlinks=true,
    linkcolor=MidnightBlue,
    urlcolor=MidnightBlue,
    citecolor=MidnightBlue
}

\captionsetup{font=small, labelfont=bf, labelsep=period}
\renewcommand{\arraystretch}{1.3}

\titleformat{\section}{\Large\bfseries}{\thesection.}{0.5em}{}[\titlerule]
\titleformat{\subsection}{\large\bfseries}{\thesubsection.}{0.5em}{}

\title{\vspace{-1cm}\textbf{H1 --- Geopolitical Risk as a Cost-Push Shock}}
\date{}
\author{}

\begin{document}
\maketitle
\vspace{-1.5cm}

\begin{tcolorbox}[colback=gray!5, colframe=gray!30, arc=3pt, boxrule=1pt,
                   left=10pt, right=10pt, top=10pt, bottom=10pt]
\textbf{H1.} A positive GPR shock raises consumer prices (inflation\,$\uparrow$) and lowers industrial production (output\,$\downarrow$) --- the \emph{cost-push signature}.
\end{tcolorbox}

%% ====================================================================
\section{Data}
%% ====================================================================

\subsection{Variables}

Table~\ref{tab:vars} lists all variables. The H1 baseline LP uses only the shock and outcome variables. Robustness controls and Stage\,4+ channel variables are loaded but not included in the baseline specification.

\begin{table}[H]
\centering
\caption{Variable Definitions}
\label{tab:vars}
\small
\begin{tabularx}{\textwidth}{@{}l l l l l l X@{}}
\toprule
\textbf{Variable} & \textbf{Code} & \textbf{Source} & \textbf{Freq.} & \textbf{Start} & \textbf{Unit} & \textbf{Transformation} \\
\midrule
\multicolumn{7}{@{}l}{\textit{Shock variables}} \\
GPR (headline) & --- & C\&I\footnotemark & M & 1900 & Index & $\log(\text{GPR})$ \\
GPT (threats)  & --- & C\&I & M & 1900 & Index & $\log(\text{GPT})$ \\
GPA (acts)     & --- & C\&I & M & 1900 & Index & $\log(\text{GPA})$ \\
\addlinespace
\multicolumn{7}{@{}l}{\textit{Outcome variables}} \\
Ind.\ Production & \texttt{INDPRO}   & FRED & M & 1919 & Idx (2017=100) & $\log(\text{IP})$ \\
CPI (headline)   & \texttt{CPIAUCSL} & FRED & M & 1947 & Idx (82--84=100) & $\log(\text{CPI})$ \\
\addlinespace
\multicolumn{7}{@{}l}{\textit{Robustness controls (not in baseline)}} \\
Unemployment & \texttt{UNRATE}  & FRED & M & 1948 & Percent & Level \\
VIX          & \texttt{VIXCLS}  & FRED & D$\to$M & 1990 & Index & $\log(\text{VIX})$ \\
\bottomrule
\end{tabularx}
\end{table}
\footnotetext{Caldara \& Iacoviello, \url{https://www.matteoiacoviello.com/gpr.htm}. Constructed from newspaper article shares across 11 broadsheet papers. GPT captures \textit{threats} (forward-looking); GPA captures \textit{acts} (realised events).}

\subsection{Stationarity}

\begin{table}[H]
\centering
\caption{Expected Integration Order (verified by \textbf{ADF tests} in \texttt{s02\_explore\_data.m})}
\label{tab:adf}
\begin{tabular}{@{}l c c c@{}}
\toprule
\textbf{Variable} & \textbf{Level} & \textbf{1st Difference} & \textbf{Order} \\
\midrule
$\log(\text{GPR}),\; \log(\text{GPT}),\; \log(\text{GPA})$ & Stationary / borderline & --- & $I(0)$ \\
$\log(\text{IP}),\; \log(\text{CPI})$ & Unit root & Stationary & $I(1)$ \\
UNRATE & Unit root (persistent) & Stationary & $I(1)$ \\
$\log(\text{VIX})$ & Stationary / borderline & --- & $I(0)$ \\
\bottomrule
\end{tabular}
\end{table}

\noindent \textbf{The LP uses cumulative log differences} as the LHS ($y_{t+h} - y_{t-1}$), so \textbf{$I(1)$ variables are handled without pre-differencing}.

\subsection{Sample Period}

\begin{table}[H]
\centering
\caption{Sample Design}
\label{tab:sample}
\begin{tabular}{@{}l l p{10cm}@{}}
\toprule
\textbf{Parameter} & \textbf{Value} & \textbf{Rationale} \\
\midrule
Start & 1985:01 & Post-Volcker disinflation; single monetary-policy regime \\
End   & 2025:12 & Latest FRED vintage (April 2026) \\
$T$   & $\approx$\,492 months & \\
COVID & Dummy 2020:03--2021:12 & Robustness: pre-COVID truncation (end 2020:02) \\
\bottomrule
\end{tabular}
\end{table}

\noindent\textit{Binding constraint.} VIX (\texttt{VIXCLS}) data begins in 1990:01. Since the baseline sample starts in 1985:01, Robustness~2 (which adds VIX as a control) loses the first 5~years and uses $\approx$\,432 observations instead of~492. The baseline and Robustness~1 are unaffected.

%% ====================================================================
\section{Model Specification}
%% ====================================================================

\subsection{Local Projection (Jord\`a, 2005)}

For each horizon $h = 0, 1, \ldots, H$, we estimate by OLS:

\begin{equation}
\underbrace{y_{t+h} - y_{t-1}}_{\text{cumulative log change}}
\;=\;
\alpha_h
\;+\;
\beta_h \cdot \text{shock}_t
\;+\;
\sum_{\ell=1}^{L} \gamma_{h,\ell}\, y_{t-\ell}
\;+\;
\sum_{\ell=1}^{L} \delta_{h,\ell}\, \text{shock}_{t-\ell}
\;+\;
\varepsilon_{t+h}
\label{eq:lp}
\end{equation}

\noindent The sequence $\{\beta_0, \beta_1, \ldots, \beta_H\}$ traces out the impulse response function (IRF). Table~\ref{tab:spec} describes each component.

\begin{table}[H]
\centering
\caption{LP Specification --- Component Details}
\label{tab:spec}
\small
\begin{tabular}{@{}l l p{8.5cm}@{}}
\toprule
\textbf{Component} & \textbf{Notation} & \textbf{Description} \\
\midrule
Dependent variable & $y_{t+h} - y_{t-1}$ & Cumulative log change in the outcome from $t{-}1$ to $t{+}h$. Outcome is $\log(\text{IP})$ or $\log(\text{CPI})$. Using cumulative differences (rather than $h$-step-ahead levels) ensures the LP is robust to $I(1)$ variables. \\
\addlinespace
Shock & $\text{shock}_t$ & Contemporaneous value of $\log(\text{GPR})$, $\log(\text{GPT})$, or $\log(\text{GPA})$. Three parallel regressions, one per shock variant. \\
\addlinespace
Own lags of outcome & $y_{t-1}, \ldots, y_{t-L}$ & Absorb persistence and serial correlation in the outcome. Prevent spurious long-horizon correlations that would arise if the outcome's autoregressive structure were omitted. \\
\addlinespace
Own lags of shock & $\text{shock}_{t-1}, \ldots, \text{shock}_{t-L}$ & Absorb persistence in the GPR series. Control for dynamic feedback: if past GPR predicts both current GPR and future macro outcomes, omitting shock lags would bias $\beta_h$ upward. \\
\addlinespace
Constant & $\alpha_h$ & Horizon-specific intercept. \\
\addlinespace
$\beta_h$ (key) & --- & The impulse response at horizon $h$. Measures the cumulative log-point change in the outcome $h$ months after a unit increase in $\log(\text{shock})$. \\
\bottomrule
\end{tabular}
\end{table}

\subsection{Why No Additional Controls in the Baseline?}

The H1 baseline deliberately excludes variables like oil prices, the federal funds rate, and credit spreads. These are \textbf{mediators} (transmission channels through which GPR affects the outcome), not confounders. Controlling for a mediator removes part of the causal effect we are trying to measure, attenuating $\beta_h$.

\begin{table}[H]
\centering
\caption{Mediator vs.\ Confounder Classification}
\label{tab:mediator}
\small
\begin{tabular}{@{}l l l p{5.5cm}@{}}
\toprule
\textbf{Variable} & \textbf{Type} & \textbf{Causal path} & \textbf{Why excluded from baseline} \\
\midrule
Oil (WTI)     & Mediator & GPR $\to$ oil $\to$ CPI          & Removes the commodity-price channel --- the dominant supply-side mechanism \\
FFR           & Mediator & GPR $\to$ $\pi$ $\to$ Fed $\to$ FFR & Removes the monetary-policy response channel \\
Credit spread & Mediator & GPR $\to$ risk premia $\to$ spreads & Removes the financial channel \\
\midrule
UNRATE        & Pot.\ confounder & See Robustness~1 & Included in robustness only \\
VIX           & Pot.\ confounder & See Robustness~2 & Included in robustness only \\
\bottomrule
\end{tabular}
\end{table}

\noindent\textbf{Identification assumption} (Caldara \& Iacoviello, 2022): GPR is driven by geopolitical events (wars, terrorism, political crises) that are exogenous to the US business cycle within the month. Under this assumption, own lags are sufficient for consistent estimation of~$\beta_h$.

%% ====================================================================
\section{Robustness Controls}
%% ====================================================================

Two robustness specifications progressively \textbf{add potential confounders to test whether $\beta_h$ is sensitive to omitted-variable bias}. The LP becomes:

\begin{equation}
y_{t+h} - y_{t-1}
\;=\;
\alpha_h + \beta_h \cdot \text{shock}_t
+ \sum_{\ell=1}^{L} \gamma_{h,\ell}\, y_{t-\ell}
+ \sum_{\ell=1}^{L} \delta_{h,\ell}\, \text{shock}_{t-\ell}
+ \sum_{\ell=1}^{L} \boldsymbol{\phi}_{h,\ell}'\, \mathbf{W}_{t-\ell}
+ \varepsilon_{t+h}
\label{eq:lp_robust}
\end{equation}

\noindent where $\mathbf{W}$ contains the robustness controls. Table~\ref{tab:robust} details each specification.

\begin{table}[H]
\centering
\caption{Robustness Specifications and Justification}
\label{tab:robust}
\small
\begin{tabular}{@{}l l p{10cm}@{}}
\toprule
\textbf{Spec.} & $\mathbf{W}$ & \textbf{Justification} \\
\midrule
Baseline & (none) & Pure total effect. Assumes GPR exogeneity (C\&I, 2022). \\
\addlinespace
Robust.\,1 & UNRATE &
\textbf{Why confounder:} The unemployment rate proxies the business-cycle state. A recession may (a)~increase geopolitical tensions through political instability, raising GPR, and (b)~independently affect inflation via the Phillips curve. Without controlling for the cycle, $\beta_h$ could partly reflect business-cycle comovement rather than the causal GPR effect. \par
\textbf{Why not mediator:} The primary GPR transmission channels are supply-side (oil, commodities, expectations). GPR does not primarily affect inflation \textit{through} unemployment. \par
\textbf{Literature:} Caldara et al.\ (2026, JIE) include output/employment controls in their country-panel. Ramey (2016, \textit{Handbook}) recommends controlling for the business cycle when the shock's exogeneity is imperfect. \\
\addlinespace
Robust.\,2 & \parbox[t]{1.8cm}{UNRATE,\\$\log(\text{VIX})$} &
\textbf{Why confounder:} Global uncertainty (VIX) could simultaneously raise news attention to geopolitical events (inflating GPR measurement) and depress investment/consumption (affecting macro outcomes). VIX measures \textit{financial-market} uncertainty --- a different dimension from GPR's \textit{geopolitical} uncertainty. \par
\textbf{Why not mediator:} If GPR directly caused VIX (GPR $\to$ VIX $\to$ output), then VIX would be a mediator. Caldara \& Iacoviello (2022, Table~3) show GPR and VIX are moderately correlated ($\rho \approx 0.3$--$0.4$) but capture distinct constructs. The shared variation likely reflects common drivers (global crises), not unidirectional causation. \par
\textbf{Literature:} Baker, Bloom \& Davis (2016, QJE) show EPU/VIX and GPR are related but not equivalent uncertainty measures. \par
\textbf{Sample note:} VIX starts 1990:01, so this specification uses $\approx$\,432 months (vs.\ 492 baseline). \\
\bottomrule
\end{tabular}
\end{table}

\noindent\textbf{Interpretation.} If $\beta_h$ remains stable across all three specifications, the GPR exogeneity assumption is supported. If $\beta_h$ changes substantially when adding UNRATE or VIX, this signals potential confounding that must be discussed.

%% ====================================================================
\section{Lag Selection}
%% ====================================================================

\subsection{Baseline Choice: $L = 12$}

\begin{table}[H]
\centering
\caption{Rationale for $L = 12$}
\label{tab:lag_rationale}
\begin{tabularx}{\textwidth}{@{}l X@{}}
\toprule
\textbf{Criterion} & \textbf{Detail} \\
\midrule
Literature precedent & Caldara \& Iacoviello (2022, AER) use 12 lags in their monthly VAR. Our LP uses the same lag length for direct comparability. \\
Seasonality absorption & 12 monthly lags span one full year, absorbing seasonal autocorrelation in CPI and IP without requiring seasonal dummies. \\
Team consistency & The team's preliminary VAR work (\texttt{Prelim\_OLS\_MICH}) uses $L = 12$. \\
Degrees of freedom & With $T \approx 492$, $L = 12$, and 2 variables lagged (shock + outcome): $2 \times 12 + 2 = 26$ RHS parameters. At $h = 36$: $T_{\text{eff}} \approx 444$, well above any concern. \\
\bottomrule
\end{tabularx}
\end{table}

\subsection{Cross-Check: VAR-Based AIC/BIC}

To verify that $L = 12$ is not an arbitrary choice, we estimate a reduced-form VAR($p$) on the system $[\text{shock},\, \log(\text{IP}),\, \log(\text{CPI})]$ for $p = 1, \ldots, 24$ and compute:

\begin{align}
\text{AIC}(p) &= \log|\hat{\Sigma}_p| + \frac{2K^2 p}{T_{\text{eff}}} \label{eq:aic} \\[4pt]
\text{BIC}(p) &= \log|\hat{\Sigma}_p| + \frac{K^2 p \log T_{\text{eff}}}{T_{\text{eff}}} \label{eq:bic}
\end{align}

\noindent where $\hat{\Sigma}_p$ is the MLE residual covariance of the VAR($p$), $K = 3$ variables, and $T_{\text{eff}} = T - p$. This is implemented in \texttt{lp\_lag\_select.m} and run at the start of \texttt{s03\_run\_h1.m}.

\begin{table}[H]
\centering
\caption{Interpretation of AIC/BIC Results}
\label{tab:lag_interpret}
\begin{tabularx}{\textwidth}{@{}l X@{}}
\toprule
\textbf{Outcome} & \textbf{Action} \\
\midrule
AIC and/or BIC select $p = 12$ & $L = 12$ is confirmed by data-driven criteria. \\
BIC selects $p < 12$, AIC selects $p \geq 12$ & BIC penalises parameters more heavily; $L = 12$ is conservative but defensible. Report BIC-optimal lag as robustness. \\
Both select $p \ll 12$ (e.g., $p \leq 6$) & $L = 12$ may be over-parameterised. Report results for both $L = 12$ and $L = p^*_{\text{BIC}}$ and compare. \\
Both select $p > 12$ & $L = 12$ may be insufficient. Extend to $L = p^*_{\text{AIC}}$ as robustness. \\
\bottomrule
\end{tabularx}
\end{table}

\noindent\textit{Note.} LP does not require VAR-based lag selection, since LP and VAR use lags differently: in VAR, lags define the system dynamics; in LP, lags are controls for conditional exogeneity. The VAR AIC/BIC cross-check serves as a \textit{comparability check} with the existing literature, not as a model-selection criterion for the LP itself. See L\"utkepohl (2005, Ch.~4.3) and Plagborg-M\o ller \& Wolf (2021, \textit{Econometrica}).

%% ====================================================================
\section{Parameters}
%% ====================================================================

\begin{table}[H]
\centering
\caption{LP Estimation Parameters}
\label{tab:params}
\begin{tabular}{@{}l l p{8.5cm}@{}}
\toprule
\textbf{Parameter} & \textbf{Value} & \textbf{Rationale} \\
\midrule
$H$ (max horizon) & 36 months & 3 years; captures hump-shaped response documented in Caldara et al.\ (2024); avoids degree-of-freedom loss \\
$L$ (lag length)   & 12 & See Section~4; cross-checked by VAR AIC/BIC \\
$\alpha$ (CI level) & 0.10 & 90\% confidence bands (standard in macro LP literature) \\
\midrule
\multicolumn{3}{@{}l}{\textit{Inference}} \\
SE method & Newey--West HAC & Bandwidth $= \max\!\bigl(h,\; \lfloor 4(T/100)^{2/9} \rfloor\bigr)$ \\
Critical value & $z_{0.05} = 1.645$ & Normal approximation (90\% two-sided) \\
\bottomrule
\end{tabular}
\end{table}

\noindent LP residuals are $\text{MA}(h)$ by construction (Jord\`a, 2005), so HAC correction is essential for valid inference at all horizons $h > 0$.

%% ====================================================================
\section{Deliverable --- 6-Panel IRF Figure}
%% ====================================================================

\begin{table}[H]
\centering
\caption{IRF Panel Layout (3 shocks $\times$ 2 outcomes)}
\label{tab:panels}
\begin{tabular}{@{}l c c@{}}
\toprule
 & \textbf{IP response} & \textbf{CPI response} \\
 & (expected: $-$) & (expected: $+$) \\
\midrule
GPR (headline) & $\beta_h\!\bigl(\log\text{IP} \mid \text{GPR}\bigr)$ & $\beta_h\!\bigl(\log\text{CPI} \mid \text{GPR}\bigr)$ \\
GPT (threats)  & $\beta_h\!\bigl(\log\text{IP} \mid \text{GPT}\bigr)$ & $\beta_h\!\bigl(\log\text{CPI} \mid \text{GPT}\bigr)$ \\
GPA (acts)     & $\beta_h\!\bigl(\log\text{IP} \mid \text{GPA}\bigr)$ & $\beta_h\!\bigl(\log\text{CPI} \mid \text{GPA}\bigr)$ \\
\bottomrule
\end{tabular}
\end{table}

\noindent Each panel: IRF line with shaded 90\% Newey--West HAC confidence band; $x$-axis: horizon 0--36 months. Produced for all three specifications (baseline, robustness\,1, robustness\,2) plus an overlay comparison figure.

%% ====================================================================
\section{Success Criterion}
%% ====================================================================

H1 is supported if, for at least one shock variant (GPR / GPT / GPA) at horizons $h \approx 6$--24 months:

\begin{itemize}[itemsep=4pt, leftmargin=1.5em]
    \item CPI impulse response is \textbf{significantly positive} (inflation\,$\uparrow$)
    \item IP impulse response is \textbf{significantly negative} (output\,$\downarrow$)
\end{itemize}

\noindent This constitutes the \textbf{cost-push signature}. The result should be robust across specifications (Table~\ref{tab:robust}).

%% ====================================================================
\section{Code Structure}
%% ====================================================================

\subsection{Folder Layout}

\begin{verbatim}
Feedback_Pf.Baek/H1/
├── run_all.m                  <- master script (run this one file)
├── README_H1.md / .tex        this document
├── scripts/
│   ├── download_data.sh       Step 0: download raw data (bash)
│   ├── s01_load_data.m        Step 1: load, merge, transform
│   ├── s02_explore_data.m     Step 2: descriptive analysis
│   └── s03_run_h1.m           Step 3: LP estimation + robustness
├── code/                      reusable functions (called by s03)
│   ├── lp_estimate.m          LP estimation + Newey-West SE
│   ├── lp_newey_west.m        Newey-West HAC standard errors
│   ├── lp_lag_select.m        VAR-based AIC/BIC lag selection
│   └── lp_plot_irf.m          IRF plotting with confidence bands
├── data/
│   ├── raw/                   populated by download_data.sh
│   └── h1_baseline.mat        created by s01 (not committed)
└── output/                    figures + results (created by scripts)
\end{verbatim}

\subsection{How to Run}

\textbf{Prerequisites:} MATLAB (R2020a or later recommended). Git Bash or WSL for the download script.

\medskip
\noindent\textbf{Step 0 --- Download raw data} (terminal, one-time):

\begin{verbatim}
cd Feedback_Pf.Baek/H1
bash scripts/download_data.sh
\end{verbatim}

\noindent This fetches FRED series (INDPRO, CPIAUCSL, UNRATE, VIXCLS, FEDFUNDS, DCOILWTICO) and the Caldara--Iacoviello GPR index into \texttt{data/raw/}. Run once, or again to refresh data.

\medskip
\noindent\textbf{Step 1 --- Run the full pipeline} (MATLAB, one command):

\begin{verbatim}
run('Feedback_Pf.Baek/H1/run_all.m')
\end{verbatim}

\noindent \texttt{run\_all.m} executes the entire pipeline in order:

\begin{table}[H]
\centering
\caption{Pipeline Execution Order}
\label{tab:execution}
\small
\begin{tabularx}{\textwidth}{@{}c l X l@{}}
\toprule
\textbf{Step} & \textbf{Script} & \textbf{What it does} & \textbf{Output} \\
\midrule
1/3 & \texttt{s01\_load\_data.m} & Load raw CSV/XLS $\to$ daily-to-monthly $\to$ merge $\to$ log transform $\to$ save & \texttt{h1\_baseline.mat} \\
2/3 & \texttt{s02\_explore\_data.m} & Summary stats, time series plots, ADF tests, correlations, ACF & \texttt{output/fig\_*.png} \\
3/3 & \texttt{s03\_run\_h1.m} & Lag selection $\to$ Baseline LP $\to$ Robustness 1\,\&\,2 $\to$ comparison & \texttt{output/fig\_h1\_*.png} \\
\bottomrule
\end{tabularx}
\end{table}

\noindent Alternatively, run each step individually:

\begin{verbatim}
run('Feedback_Pf.Baek/H1/scripts/s01_load_data.m')    % step 1
run('Feedback_Pf.Baek/H1/scripts/s02_explore_data.m')  % step 2
run('Feedback_Pf.Baek/H1/scripts/s03_run_h1.m')        % step 3
\end{verbatim}

\noindent Functions in \texttt{code/} (\texttt{lp\_estimate.m}, \texttt{lp\_newey\_west.m}, \texttt{lp\_lag\_select.m}, \texttt{lp\_plot\_irf.m}) are called internally by \texttt{s03} --- do not run them directly.

%% ====================================================================
\section*{References}
%% ====================================================================

\begin{itemize}[itemsep=4pt, leftmargin=1.5em]
    \item Baker, S.R., Bloom, N. \& Davis, S.J. (2016). ``Measuring Economic Policy Uncertainty.'' \textit{Quarterly Journal of Economics}, 131(4), 1593--1636.
    \item Caldara, D. \& Iacoviello, M. (2022). ``Measuring Geopolitical Risk.'' \textit{American Economic Review}, 112(4), 1194--1225.
    \item Caldara, D., Conlisk, S., Iacoviello, M. \& Penn, M. (2026). ``Do Geopolitical Risks Raise or Lower Inflation?'' \textit{Journal of International Economics}, 159.
    \item Jord\`a, \`O. (2005). ``Estimation and Inference of Impulse Responses by Local Projections.'' \textit{American Economic Review}, 95(1), 161--182.
    \item L\"utkepohl, H. (2005). \textit{New Introduction to Multiple Time Series Analysis}. Springer.
    \item Newey, W.K. \& West, K.D. (1987). ``A Simple, Positive Semi-Definite, Heteroskedasticity and Autocorrelation Consistent Covariance Matrix.'' \textit{Econometrica}, 55(3), 703--708.
    \item Plagborg-M\o ller, M. \& Wolf, C.K. (2021). ``Local Projections and VARs Estimate the Same Impulse Responses.'' \textit{Econometrica}, 89(2), 955--980.
    \item Ramey, V.A. (2016). ``Macroeconomic Shocks and Their Propagation.'' \textit{Handbook of Macroeconomics}, Vol.~2, 71--162.
\end{itemize}

\end{document}
